\documentclass[10pt,a4paper]{amsart}
\usepackage{amsmath,amssymb,amsthm,graphicx}
% Retained minimal preamble for the Stacks Project chapter more-algebra.tex (More on Algebra).
% The chapter file pinned-original-source.tex carries only \input{preamble} on its first line, so no
% preamble text lives inside the chapter itself.  The Stacks Project uses one single shared preamble for
% every chapter; the snapshot stages it identically as sources/stacks-chow/preamble.tex and
% sources/stacks-etale-cohomology/preamble.tex.  The block below is copied mechanically, line for line,
% from sources/stacks-chow/preamble.tex lines 14-18 (the Xy-pic setup) and 154-217 (the native theorem
% declarations together with the mathematical macros).  Nothing was added, removed or redefined; the
% wrapper class and the amsmath/amssymb/amsthm/graphicx base packages are supplied by the assembler.
% Chapter: more-algebra.tex -- More on Algebra.
% For commutative diagrams we use Xy-pic
\usepackage[all]{xy}

% We use 2cell for 2-commutative diagrams.
\xyoption{2cell}

% Theorem environments.
%
\theoremstyle{plain}
\newtheorem{theorem}[subsection]{Theorem}
\newtheorem{proposition}[subsection]{Proposition}
\newtheorem{lemma}[subsection]{Lemma}

\theoremstyle{definition}
\newtheorem{definition}[subsection]{Definition}
\newtheorem{example}[subsection]{Example}
\newtheorem{exercise}[subsection]{Exercise}
\newtheorem{situation}[subsection]{Situation}

\theoremstyle{remark}
\newtheorem{remark}[subsection]{Remark}
\newtheorem{remarks}[subsection]{Remarks}

\numberwithin{equation}{subsection}

% Macros
%
\def\lim{\mathop{\mathrm{lim}}\nolimits}
\def\colim{\mathop{\mathrm{colim}}\nolimits}
\def\Spec{\mathop{\mathrm{Spec}}}
\def\Hom{\mathop{\mathrm{Hom}}\nolimits}
\def\Ext{\mathop{\mathrm{Ext}}\nolimits}
\def\SheafHom{\mathop{\mathcal{H}\!\mathit{om}}\nolimits}
\def\SheafExt{\mathop{\mathcal{E}\!\mathit{xt}}\nolimits}
\def\Sch{\mathit{Sch}}
\def\Mor{\mathop{\mathrm{Mor}}\nolimits}
\def\Ob{\mathop{\mathrm{Ob}}\nolimits}
\def\Sh{\mathop{\mathit{Sh}}\nolimits}
\def\NL{\mathop{N\!L}\nolimits}
\def\CH{\mathop{\mathrm{CH}}\nolimits}
\def\proetale{{pro\text{-}\acute{e}tale}}
\def\etale{{\acute{e}tale}}
\def\QCoh{\mathit{QCoh}}
\def\Ker{\mathop{\mathrm{Ker}}}
\def\Im{\mathop{\mathrm{Im}}}
\def\Coker{\mathop{\mathrm{Coker}}}
\def\Coim{\mathop{\mathrm{Coim}}}

% Boxtimes
%
\DeclareMathSymbol{\boxtimes}{\mathbin}{AMSa}{"02}

%
% Macros for moduli stacks/spaces
%
\def\QCohstack{\mathcal{QC}\!\mathit{oh}}
\def\Cohstack{\mathcal{C}\!\mathit{oh}}
\def\Spacesstack{\mathcal{S}\!\mathit{paces}}
\def\Quotfunctor{\mathrm{Quot}}
\def\Hilbfunctor{\mathrm{Hilb}}
\def\Curvesstack{\mathcal{C}\!\mathit{urves}}
\def\Polarizedstack{\mathcal{P}\!\mathit{olarized}}
\def\Complexesstack{\mathcal{C}\!\mathit{omplexes}}
% \Pic is the operator that assigns to X its picard group, usage \Pic(X)
% \Picardstack_{X/B} denotes the Picard stack of X over B
% \Picardfunctor_{X/B} denotes the Picard functor of X over B
\def\Pic{\mathop{\mathrm{Pic}}\nolimits}
\def\Picardstack{\mathcal{P}\!\mathit{ic}}
\def\Picardfunctor{\mathrm{Pic}}
\def\Deformationcategory{\mathcal{D}\!\mathit{ef}}
\usepackage{cleveref}
\crefname{originallemma}{original lemma}{original lemmas}
\Crefname{originallemma}{Original lemma}{Original lemmas}
\makeatletter
\newlabel{more-algebra-lemma-module-over-fibre-product}{{07RU}{source}{}{sourcepointer-1}{}}
\newlabel{more-algebra-lemma-module-over-fibre-product@cref}{{[originallemma][0][]07RU}{[1][0][]source}{}{}{}}
\makeatother
\crefname{originalsection}{original section}{original sections}
\Crefname{originalsection}{Original section}{Original sections}
\makeatletter
\newlabel{situation-module-over-fibre-product}{{08KH}{source}{}{sourcepointer-2}{}}
\newlabel{situation-module-over-fibre-product@cref}{{[originalsection][0][]08KH}{[1][0][]source}{}{}{}}
\makeatother
\crefname{originalequation}{original equation}{original equations}
\Crefname{originalequation}{Original equation}{Original equations}
\makeatletter
\newlabel{equation-functor}{{08KI}{source}{}{sourcepointer-3}{}}
\newlabel{equation-functor@cref}{{[originalequation][0][]08KI}{[1][0][]source}{}{}{}}
\makeatother
\crefname{originallemma}{original lemma}{original lemmas}
\Crefname{originallemma}{Original lemma}{Original lemmas}
\makeatletter
\newlabel{lemma-modules}{{0D2F}{source}{}{sourcepointer-4}{}}
\newlabel{lemma-modules@cref}{{[originallemma][0][]0D2F}{[1][0][]source}{}{}{}}
\makeatother
\title{The functor from modules over a fibre product has a right adjoint given by a fibre product of modules}
\author{The Stacks Project Authors}
\date{stacks-project:more-algebra.tex:local-snapshot-20261012; GNU FDL 1.2 or later}
\begin{document}
\maketitle
\setcounter{subsection}{1}
\noindent\textit{Editorial note.} All delivered mathematical text is the author's own, copied from the pinned Stacks Project chapter file more-algebra.tex; no statement or proof step was rewritten, reordered or removed.\par\medskip
\noindent\textit{Editorial note.} The chapter file pinned-original-source.tex carries only \textbackslash{}input\{preamble\} on its first line, so no preamble text lives inside the chapter. The retained file sources/stacks-more-algebra/retained-preamble.tex is a mechanical, line-for-line copy of the project's single shared preamble as staged in the snapshot, namely sources/stacks-chow/preamble.tex lines 14-18 (Xy-pic setup) and 154-217 (theorem declarations and mathematical macros). Nothing was added, removed or redefined, no package was added, and the excerpt serves only as preamble, never as a span coordinate.\par\medskip
\noindent\textit{Editorial note.} Both counted spans start at the \textbackslash{}begin line of their environment, so the statement and the author's proof are delivered whole; surrounding prose of the section is deliberately outside the counted spans and is carried as standard prerequisites instead.\par\medskip
\noindent\textit{Editorial note.} The citation fields (source\_id, source\_version, authors, source\_url, license.name, source\_title) are identical to those of the cards already published from this chapter, so that the batch selector treats them as one source instead of deferring them as a same-source clash.\par\medskip
\noindent\textit{Allowed standard prerequisites (editorial).}\begin{itemize}
\item The Situation of a fibre product of rings with a module over it, and the functor displayed before the statement. situation environment at source lines 521-542 and the functor (equation-functor) at source lines 625-632, both before the counted statement. Fixes the data the statement refers to, namely the diagram of rings B = B' x\_A A', the module M', the descent datum phi and the functor whose right adjoint is computed.
\item The functor that sends a triple (N, M', phi) to its two base changes has a right adjoint, proved in the chapter's module lemma. the author cites this at proof lines 659-660 as Lemma lemma-modules (tag 0D2F), a statement of the same chapter. Supplies the adjointness half of the assertion, which the author does not reprove; only the identity statement is carried out inside the counted span.
\item Exactness and flatness-free manipulation of the fibre product of modules, in particular the identification of kernels of the horizontal maps. used at proof lines 661-679, where the kernels I, J, K, L of the two diagrams are introduced and compared. Supplies the objects in terms of which all twelve observations are formulated.
\end{itemize}\medskip
\ifdefined\hypertarget\hypertarget{sourcepointer-1}{}\fi
\noindent\textit{Source reference.} Original lemma 07RU: focused-seeds/AGP/sources/stacks-snapshot-20261011/tags.txt, source lines 9879--9879. The lemma's own canonical tag, quoted so that readers can resolve the label inside the delivered text; it points at the counted statement itself and carries no proof content.\par\medskip
\ifdefined\hypertarget\hypertarget{sourcepointer-2}{}\fi
\noindent\textit{Source reference.} Original section 08KH: focused-seeds/AGP/sources/stacks-snapshot-20261011/tags.txt, source lines 525--525. Navigation pointer to the Situation that fixes the rings, the modules and the comparison map of the statement; it contains no proof content and the counted argument does not use anything from it beyond the naming of the data.\par\medskip
\ifdefined\hypertarget\hypertarget{sourcepointer-3}{}\fi
\noindent\textit{Source reference.} Original equation 08KI: focused-seeds/AGP/sources/stacks-snapshot-20261011/tags.txt, source lines 629--629. Navigation pointer to the displayed functor whose right adjoint the statement describes; quoting the display is a matter of notation, and the identification of the adjoint and the identity statement are carried out inside the counted proof span.\par\medskip
\ifdefined\hypertarget\hypertarget{sourcepointer-4}{}\fi
\noindent\textit{Source reference.} Original lemma 0D2F: focused-seeds/AGP/sources/stacks-snapshot-20261011/tags.txt, source lines 16380--16380. Navigation pointer to the chapter's module lemma that the author cites for the adjointness half of the statement; the second half, which is the substance of the card, is proved inside the counted span.\par\medskip
\begin{lemma}
\label{lemma-module-over-fibre-product}
In Situation \ref{situation-module-over-fibre-product}
the functor (\ref{equation-functor}) has a right adjoint, namely
the functor
$$
F : (N, M', \varphi) \longmapsto N \times_{\varphi, M} M'
$$
where $M = M'/IM'$. Moreover, the composition of $F$ with
(\ref{equation-functor}) is the identity functor on
$\text{Mod}_B \times_{\text{Mod}_A} \text{Mod}_{A'}$. In other words,
setting $N' = N \times_{\varphi, M} M'$ we have
$N' \otimes_{B'} B = N$ and $N' \otimes_{B'} A' = M'$.
\end{lemma}

\begin{proof}
The adjointness statement follows from the more general
Lemma \ref{lemma-modules}.
To prove the final assertion, recall that
$B' = B \times_A A'$ and $N' = N \times_{\varphi, M} M'$ and extend
these equalities to 
$$
\vcenter{
\xymatrix{
A & A' \ar[l] & I \ar[l] \\
B \ar[u] & B' \ar[l] \ar[u] & J \ar[l] \ar[u]
}
}
\quad\text{and}\quad
\vcenter{
\xymatrix{
M & M' \ar[l] & K \ar[l] \\
N \ar[u]_\varphi & N' \ar[l] \ar[u] & L \ar[l] \ar[u]
}
}
$$
where $I, J, K, L$ are the kernels of the horizontal maps of the original
diagrams. We present the proof as a sequence of observations:
\begin{enumerate}
\item $K = IM'$ (see statement lemma),
\item $B' \to B$ is surjective with kernel $J$ and $J \to I$ is bijective,
\item $N' \to N$ is surjective with kernel $L$ and $L \to K$ is bijective,
\item $JN' \subset L$,
\item $\Im(N \to M)$ generates $M$ as an $A$-module
(because $N \otimes_B A = M$),
\item $\Im(N' \to M')$ generates $M'$ as an $A'$-module
(because it holds modulo $K$ and $L$ maps isomorphically to $K$),
\item $JN' = L$ (because $L \cong K = I M'$
is generated by images of elements $x n'$ with $x \in I$ and
$n' \in N'$ by the previous statement),
\item $N' \otimes_{B'} B = N$ (because $N = N'/L$, $B = B'/J$, and
the previous statement),
\item there is a map $\gamma : N' \otimes_{B'} A' \to M'$,
\item $\gamma$ is surjective (see above),
\item the kernel of the composition $N' \otimes_{B'} A' \to M' \to M$ 
is generated by elements $l \otimes 1$ and $n' \otimes x$ with
$l \in K$, $n' \in N'$, $x \in I$ (because $M = N \otimes_B A$ by assumption
and because $N' \to N$ and $A' \to A$ are surjective with kernels
$L$ and $I$),
\item any element of $N' \otimes_{B'} A'$ in the submodule generated
by the elements $l \otimes 1$ and $n' \otimes x$ with
$l \in L$, $n' \in N'$, $x \in I$ can be written as $l \otimes 1$
for some $l \in L$ (because $J$ maps isomorphically to $I$ we see
that $n' \otimes x = n'x \otimes 1$ in $N' \otimes_{B'} A'$;
similarly $x n' \otimes a' = n' \otimes xa' = n'(xa') \otimes 1$
in $N' \otimes_{B'} A'$ when $n' \in N'$, $x \in J$ and $a' \in A'$;
since we have seen that $JN' = L$ this proves the assertion),
\item the kernel of $\gamma$ is zero (because by (10) and (11) any element of
the kernel is of the form $l \otimes 1$ with $l \in L$ which
is mapped to $l \in K \subset M'$ by $\gamma$).
\end{enumerate}
This finishes the proof.
\end{proof}

\end{document}
